\section{Experimental Setup}
\label{sec:experimental-setup}

This section documents the executable benchmark used for the sequential multi-truck experiments reported in Table~2. Throughout, we use the current repository configuration as the source of truth for environment construction, training, and evaluation. Where manuscript-level conventions may differ from the executable setup, we state the discrepancy explicitly rather than smoothing it over.

\subsection{Experimental network and simulator parameterization}
\label{sec:experimental-network}

The benchmark is built on a directed California transportation network encoded in three precomputed data files: a shortest-path energy dictionary, a shortest-path travel-time dictionary, and a charging-station metadata table. In the executable graph, these files produce a network with 258 nodes, 66{,}049 directed edges, 25 charging nodes, and 209 total charging ports. Figure~6 should be introduced here as the experimental geography, showing the charger locations, candidate delivery nodes, transportation links, and the spatial variation in average energy requirements.

Delivery locations are not fixed customer nodes stored separately in the dataset. Instead, candidate delivery nodes are defined procedurally as all non-charger nodes with nonzero out-degree. Depot/start nodes are sampled from the same valid non-charger pool, so each truck begins an episode from a sampled network node that is feasible as an origin. For each truck, the delivery generator then samples an ordered sequence whose consecutive hop energies fall in the configured range of 100--250~kWh. The sampled sequence is feasibility-validated against the battery and charging network, and if a feasible sequence cannot be found under the current hop bounds, the bounds are progressively reduced until a valid sequence is obtained.

Both energy and travel time are treated as precomputed shortest-path quantities on the directed graph. The resulting network is strongly asymmetric: among the 32{,}896 bidirectional node pairs present in the executable graph, all are energy-asymmetric and 32{,}885 are travel-time asymmetric. This asymmetry is important because it prevents the benchmark from collapsing into a symmetric stylized routing problem.

The main sequential benchmark uses homogeneous electric trucks with a 400~kWh battery and a base speed of 40~km/h. Initial battery is set to full charge. Charging decisions are discretized into 12 durations, from 1 to 12 hours. Realistic charging is enabled, so DC fast charging follows a nonlinear CCCV curve rather than a constant-rate approximation. The infrastructure itself is fixed across instance families: all settings use the same 25 charging nodes and 209 aggregate charging ports, while the number of depots varies with the number of trucks because each truck receives its own sampled start node.

One implementation detail deserves an explicit cautionary note. The raw station metadata file contains 15 Level~2 sites and 10 DC fast sites, but the current graph loader maps all charger types to \texttt{DCFast} when building the executable environment. The manuscript should therefore either describe the executed setup exactly as an effectively all-fast-charging benchmark, or this preprocessing choice should be reconciled before claiming mixed charger types in the experimental section.

Overall, the network is realistic enough to induce asymmetric travel and energy costs, while still being controllable enough for systematic benchmarking across problem sizes.

\subsection{Scenario generation and problem settings}
\label{sec:scenario-generation}

We index benchmark settings by labels of the form \texttt{T}--\texttt{S}, where \texttt{T} denotes the number of trucks and \texttt{S} denotes the number of delivery stops assigned to each truck. The main benchmark reported in Table~2 uses the sequential fixed-order setting with \texttt{S}=3, namely \texttt{1T3S}, \texttt{5T3S}, \texttt{10T3S}, \texttt{30T3S}, \texttt{50T3S}, and \texttt{100T3S}. In this benchmark, delivery order is fixed in advance (\texttt{enable\_flexible\_delivery\_order = false}), so the routing decision is not which customer to serve next, but when to detour for charging and how long to charge.

Episodes are generated online at reset time. Each truck is assigned a random depot/start node and a feasible ordered delivery sequence drawn under the 100--250~kWh hop-energy constraint described above. The main sequential configuration keeps the number of stops fixed for each target setting (\texttt{allow\_variable\_num\_stops = false}) and initializes all trucks at full battery. Charger capacity is inherited from the fixed network metadata, so all settings share the same 25 charging locations and 209 total ports.

Each RL model is trained separately for its target setting; the saved checkpoint layout and runner scripts are setting-specific rather than shared across sizes. For final evaluation, the current parallel evaluator uses 100 random test scenarios per setting with deterministic seeding (\texttt{seed + episode\_idx}) so that methods are compared on matched scenarios. The repository does not expose a persistent train/validation/test split artifact; instead, training is on-policy with periodic evaluation during training. Accordingly, we do not claim an explicit split in this section unless such a split is documented separately in archived experiment logs.

If the final manuscript table is based on 200 test scenarios rather than 100, that number should be treated as a manual override sourced from the archived experiment runs rather than from the current evaluator. The executable code in this repository currently supports the 100-scenario protocol.

\subsection{Uncertainty and charging model}
\label{sec:uncertainty-model}

The realism of the benchmark comes from combining a directed asymmetric network with exogenous travel and energy uncertainty, endogenous charger contention, and nonlinear charging dynamics; see also Figure~3. Travel time is modeled as a Gaussian perturbation around the nominal shortest-path time, with a base standard-deviation factor of 0.15, a rush-hour variance multiplier of 2.0, and a standard-deviation cap of 1.0~hour. Rush-hour exposure is computed over the travel interval itself, and realized travel times are clipped to the range $[0.85, 2.5]$ times the nominal value.

Energy consumption is stochastic and correlated with traffic conditions. The base energy on a path is perturbed with a standard-deviation factor of 0.05 and clipped to the range $[0.90, 1.20]$ times nominal. The feasibility mask uses the worst-case multiplier of 1.20 when checking whether a routing action is safe, so the policy reasons about conservative energy feasibility even though realized energy remains stochastic.

Charging contention is endogenous. Each charging station has finite capacity, and arriving trucks enter a strict first-come, first-served waitlist when all ports are occupied. Waiting trucks are reactivated through event-driven wake-up logic as soon as a port becomes available, so queueing delays depend on the joint fleet trajectory rather than on an exogenous delay model.

Charging itself is nonlinear in the main benchmark. The environment uses a CCCV charging model with numerical integration time step $\Delta t = 0.01$~h. Power ramps up at low state of charge, remains at peak power through the plateau region, and tapers once the battery exceeds 80\% SoC. This is materially different from linear charging approximations that assume constant charging power throughout the session. The simulator also supports stochastic unloading times, but this mechanism is disabled in the main sequential configuration (\texttt{enable\_stochastic\_unloading = false}) and therefore should not be described as active in the benchmark unless the reported runs use a different configuration.

\subsection{Compared methods}
\label{sec:compared-methods}

We compare five method classes in the main sequential benchmark.

\paragraph{Math. Opt.}
The optimization baseline is a per-truck Gurobi mixed-integer planner that uses deterministic travel times, safety-buffered energy consumption, and no charger contention. In the evaluation scripts, this label corresponds to the best available optimization-based reference for a given setting: the standard deterministic MILP on smaller instances and the simplified robust variant on larger instances.

\paragraph{Heuristic.}
The heuristic baseline is an environment-aware greedy policy. It checks battery sufficiency, charger reachability, and post-delivery feasibility, then chooses either the next delivery or a reachable charger together with the minimum required charging duration.

\paragraph{PPO.}
This baseline uses a generic Stable-Baselines3 PPO policy with an MLP observation encoder and no graph structure.

\paragraph{MaskPPO.}
MaskPPO uses the same generic PPO family but with the environment's feasibility mask exposed during action selection. It therefore benefits from explicit invalid-action suppression without using the full graph-structured state/action representation of the proposed method.

\paragraph{GraphPPO.}
GraphPPO is the proposed heterogeneous-graph PPO policy. It operates on graph-structured observations and graph-structured feasible action sets, enabling explicit reasoning over trucks, deliveries, chargers, and their relational attributes.

Classical VRP heuristics such as Clarke--Wright savings and nearest-neighbor + 2-opt are implemented in the repository, but they should be discussed in this subsection only if they appear in the result table being reported. If Table~2 is restricted to Math.~Opt., Heuristic, PPO, MaskPPO, and GraphPPO, then the comparison narrative should stay aligned with those five roles.

\subsection{Training protocol and hyperparameters}
\label{sec:training-protocol}

GraphPPO is trained separately for each target setting. The saved sequential checkpoints used in the benchmark share the same core hyperparameters: hidden dimension 32, three graph layers, critic MLP width 256, learning rate $3\times 10^{-4}$, discount factor $\gamma=0.99$, GAE parameter $\lambda=0.95$, PPO clip coefficient 0.2, entropy coefficient 0.1, value-loss coefficient 0.01, gradient clipping at 0.5, 256 steps per update, five PPO epochs per update, and minibatch size 256. The saved sequential GraphPPO runs also use one training environment per run. In the benchmark runner shipped with the repository, both GraphPPO and the SB3 baselines are trained with a nominal budget of $10^6$ environment steps and periodic evaluation every 5{,}000 timesteps; if the final reported table uses archived runs with a different budget, that logged value should replace the runner default here.

For the generic PPO and MaskPPO baselines, the repository training script instantiates the default Stable-Baselines3 implementations unless overridden externally. The repo therefore supports a clean statement that PPO and MaskPPO use default SB3 hyperparameters in the current benchmark code, but exact non-default overrides should be reported only if they are available in archived experiment logs.

Table~\ref{tab:exp-hparams} summarizes the executable environment and the checkpoint-backed GraphPPO hyperparameters. Hardware and average training time per setting are intentionally left out here because they are not stored reliably in the checkpoints themselves and should be filled from W\&B or experiment logs when preparing the final manuscript.

\begin{table}[t]
\centering
\small
\begin{tabular}{@{}lll@{}}
\toprule
Block & Parameter & Value \\
\midrule
Environment & Trucks per setting & $\{1, 5, 10, 30, 50, 100\}$ \\
Environment & Stops per truck & $3$ \\
Environment & Battery capacity & $400$~kWh \\
Environment & Base speed & $40$~km/h \\
Environment & Initial battery & Full charge \\
Environment & Hop-energy range & $100$--$250$~kWh \\
Environment & Charge durations & $1,2,\dots,12$~h \\
Environment & Traffic std factor & $0.15$ \\
Environment & Energy std factor & $0.05$ \\
Environment & Energy feasibility multiplier & $1.20$ \\
Environment & Charging model & CCCV, $\Delta t = 0.01$~h \\
\midrule
GraphPPO & Hidden dimension & $32$ \\
GraphPPO & Graph layers & $3$ \\
GraphPPO & Critic MLP width & $256$ \\
GraphPPO & Learning rate & $3\times 10^{-4}$ \\
GraphPPO & Discount factor $\gamma$ & $0.99$ \\
GraphPPO & GAE $\lambda$ & $0.95$ \\
GraphPPO & PPO clip & $0.2$ \\
GraphPPO & Entropy coefficient & $0.1$ \\
GraphPPO & Value coefficient & $0.01$ \\
GraphPPO & Max grad norm & $0.5$ \\
GraphPPO & Steps per update & $256$ \\
GraphPPO & PPO epochs & $5$ \\
GraphPPO & Minibatch size & $256$ \\
GraphPPO & Training environments & $1$ \\
\midrule
SB3 baselines & Policy family & PPO / MaskablePPO \\
SB3 baselines & Hyperparameters & Default SB3 unless logged otherwise \\
\bottomrule
\end{tabular}
\caption{Executable environment settings and checkpoint-backed GraphPPO hyperparameters for the sequential T3S benchmark. Hardware and wall-clock training times should be filled from experiment logs.}
\label{tab:exp-hparams}
\end{table}

\subsection{Evaluation metrics and statistical protocol}
\label{sec:evaluation-metrics}

All reported results should be defined before the reader reaches the performance tables. The evaluator records reward, success rate, total deliveries completed, average final SoC, number of charging sessions, charging time, waiting time, routing time, unloading time, total operational time, and execution time. Unless stated otherwise, reported values are means and standard deviations over 100 matched random scenarios for each setting.

Execution time in the current evaluator is measured around the full episode rollout, including simulator stepping and policy decisions; it is therefore not a pure inference-only latency metric. This distinction matters when comparing learned policies against optimization-based baselines because the reported wall-clock reflects end-to-end evaluation cost.

The runtime evaluation scripts log raw episode reward rather than a normalized reward directly. When Table~2 uses normalized reward, that quantity should be described as a post-processing metric computed from matched per-episode rewards. In the current analysis pipeline, the normalization takes the form
\[
R_{\mathrm{norm}} = \frac{R_{\pi}}{R_{\mathrm{Math.\ Opt.}}},
\]
where both numerator and denominator are computed on the same seeded scenario. This normalization is performed after evaluation rather than inside the environment or the core evaluator. If the final manuscript uses a different normalization convention, the exact formula should be replaced here and tied to the analysis script or notebook that produced Table~2.

Finally, two paper-level consistency checks should be completed before freezing this section: first, confirm whether Table~2 should report the repo's current 100-scenario protocol or an archived 200-scenario test set; second, confirm whether the charger description should follow the executable all-\texttt{DCFast} mapping or a reconciled mixed-charger preprocessing pipeline.
